\documentclass[12pt]{article}
\usepackage[margin=1in]{geometry}
\usepackage{setspace}
\doublespacing

\begin{document}

\section*{Journal, Chapter 6, 7, 8, and 9}

For song \#1 I chose a song I have had in my playlist for some time now, 
\textit{Good Vibrations} by the Beach Boys, [Chapter 7, psychedelia].  
Of course, I first took a listen to the song before re-reading Chapter 7 
to figure out what specifics I wanted to quote... Although I was looking into 
other sources, I decided against utilizing them.
\newline

Upon a first listen, the song is busy without too much remarkable about it, 
except that it shifts between ``loud" and quiet very often. The textbook implies 
that Brian Wilson’s ``pocket symphony" reflects 1960s America’s fascination 
with pushing musical boundaries, and I tend to agree. Covach notes that 
``\emph{Good Vibrations consumed more studio time and budget allocation than any other pop single ever had before}" 
(Covach, p. 253). 

I also find myself agreeing with Covach’s assessment that the track 
``\emph{takes the listener on a `trip'... through unfamiliar musical territory}" 
(Covach, p. 253). According to the book, it had tack piano, Jew’s harp, and even sleigh bells. That part definitely captures the idea 
of mid-’60s America, when things felt like it was on the brink of some `breakthrough'.

However, I’m not entirely convinced by Covach’s suggestion that 
\textit{Good Vibrations} is primarily a ``psychedelic" record. Certainly Wilson’s production approach (with heavy overdubs and tape experiments) lines up with the psychedelic `aim' of expanding consciousness. Yet the track still has that buoyant, sunny California 
hook ``\emph{I’m pickin’ up good vibrations}" (Covach, p. 253), which, in my opinion, 
keeps it more in the realm of radio-friendly pop than some of the more overtly 
experimental work of, say, the Grateful Dead. I think it would be better known as a hybrid: a bold pop single that borrowed 
psychedelic production method.

Covach also connects \textit{Good Vibrations} to the overall American mentality 
of the period, pointing out that it ``\emph{demonstrated that rock could stand 
on its own as music and be taken seriously}" (Covach, p. 253). In a decade when The Beatles were raising the bar for album production, Brian 
Wilson took it upon himself to meet and even top those standards. That 
competitive environment was key to expanding the scope of rock music. 

\subsection*{Song 2}
Carole King’s \textit{“You’ve Got a Friend”} reflects the 1970s American culture of introspection and individualism, a hallmark of the era’s singer-songwriter movement. This style served as the \textit{“direct antithesis to Bowie and Cooper,”} focusing on \textit{“revealing their unmediated personal perspectives”} rather than theatrical personas (Covach, p. 323). The song’s production, featuring a \textit{“solo acoustic piano”} and \textit{“acoustic guitar just barely heard in the background,”} creates an \textit{“aural impression of sincerity and intimacy”} that resonated deeply with audiences transitioning into adulthood (Covach, p. 323).

The cultural relevance of the singer-songwriter movement lay in its appeal to listeners in the “Me Decade,” described as an era when individuals sought personal reflection and authenticity. King’s \textit{“You’ve Got a Friend”} exemplifies this ethos, offering a message of connection and emotional reassurance that reflected the societal yearning for stability amidst a rapidly changing world (Covach, p. 323).

\section*{References}

Covach, J. \& Flory, A. \textit{What’s That Sound? An Introduction to Rock and Its History} (5th ed.). W. W. Norton and Company.

\end{document}
